%% kapa: exam=1 round=29 year=2018 period=1 form=A subject=civil_law subject_ko=민법 seq=10 number=10
%% source: Q-Net 감정평가사 기출문제 (cst00309) artlSeq=5202808 file=extracted/1008/5202808/2018년 제29회 감정평가사 1차시험 기출문제/2018년 제29회 감정평가사 1차시험 기출문제/1교시 A형.pdf
%% conversion: kordoc 4.24.51 -> tools/gen1.py (auto); review=proofread-sonnet; flags=none
\begin{question}{10}
착오로 인한 의사표시에 관한 설명으로 옳지 않은 것은? (다툼이 있으면 판례에 따름)
\begin{choices}
  \item 제3자의 기망으로 표시상의 착오가 발생한 경우, 표의자는 사기를 이유로 의사표시를 취소할 수 있다.
  \item 착오로 인하여 표의자가 경제적인 불이익을 입지 않았다면, 법률행위 내용의 중요부분의 착오라고 할 수 없다.
  \item 표의자의 착오를 알고 상대방이 이를 이용한 경우에는 착오가 표의자의 중대한 과실로 발생하여도 취소할 수 있다.
  \item 당사자의 합의로 착오로 인한 의사표시의 취소에 관한 민법 제109조 제1항의 적용을 배제할 수 있다.
  \item 동기의 착오를 이유로 의사표시를 취소할 때 그 동기를 의사표시의 내용으로 하는 당사자의 합의까지는 필요 없다.
\end{choices}
\end{question}
\answer{1}
